\chapter{Floyd Warshall: Chiusura Transitiva}

\section{Definizione del problema}

\subsection{Istanza}
\begin{itemize}
    \item Un DAG $G=(V,E,w)$
    \begin{itemize}
        \item $V=\{1,2,\dots,n\}$ è l'insieme dei vertici, $|V|=n$
        \item $E \subseteq V^2$ è l'insieme degli archi
        \item $w:E \to \mathbb{R}^+$ stabilisce l'esistenza di un dato commino tra le coppie $(i,j)$ gerica di vertici
    \end{itemize}
    \item Il grafo è rappresentato in memoria mediante la sua matrice di adiacenza booleana $M=(m_{ij})$ di dimensione $n \times n$
    {
        \large

            $$
                m_{ij} = 
                \begin{cases}
                    1 (true) & \text{se } i=j \lor (i,j) \in E \\
                    0 (false) & \text{se } i \not = j \land (i,j) \not \in E
                \end{cases}
            $$
    }
\end{itemize}

\subsection{Soluzione}
Costruire la chiusura transitiva del grafo $G$, ovvero il grafo $G^* = (V,E^*)$ avente lo stesso insieme di vertici $V$ e un insieme di archi ampliato $E^* = \{(i,j) \in V^2 : \text{esiste un cammino da $i$ a $j$ in $G$}\}$.

L'output è rappresentato dalla matrice booleana $T=(t_{ij})$ in cui $t_{ij} = 1$ se esiste almeno un cammino da $i$ a $j$, e $0$ altrimenti.


\section{Definizione del sottoproblema $k(0 \leq k \leq n)$}
Fissando un indice $k$ per ogni coppia di vertici $(i,j)$ il sottoproblema consiste nello stabilire se esiste un cammino da $i$ a $j$ tale che eventuali vertici intermedi appartengono esclusivamente al sottoinsieme $\{1,2,\dots,k\}$(se $k=0$, l'insieme dei vertici intermedi è vuoto $\emptyset$)

\section{Definizione del coefficiente $\mathbf{t^{(k)_{ij}}}$}
Per ogni $k$ e per ogni coppia $(i,j)$, introduciamo la variabile booleana

{
    \large

    $$
    t^{k}_{ij} = 
    \begin{cases}
        1 (true) & \text{se esiste un cammino da $i$ a $j$ con vertici intermedi in $\{1,\dots,k\}$} \\
        0 (false) & \text{altrimenti}
    \end{cases}
    $$
}

\newpage
\section{Formulazione ricorsiva}
\subsection{Caso base $\mathbf{k=0}$}
Non essendoci vertici intermedi ammessi $\emptyset$, un cammino da $i$ a $j$ esiste se e solo se $i=j$ oppure esiste l'arco diretto $(i,j) \in E$.

I coefficienti iniziali coincidono con la matrice di adiacenza del grafo $(T^{(0)}=M)$

{
    \large

    $$
        t^{(0)}_{ij} = m_{ij}
    $$
}

\subsection{Passo ricorsivo $\mathbf{k>0}$}


{
    \large
    $$
    t_{ij}^{(k)} = t_{ij}^{k-1} \lor (t_{ik}^{(k-1)} \land t_{kj}^{(k-1)})
    $$
}

\subsubsection{Spiegazione}
Assumendo di aver calcolato i coefficienti al passo $k-1$, per determinare $t^{(k)}_{ij}$ si valutano due casi rispetto al vertice $k$

\begin{itemize}
    \item \textbf{Il vertice $\mathbf{k}$ non è un vertice intermedio del cammino}\\
    Il cammina usa solo vertici in $\{1,\dots,k-1\}$, quindi il valore è $t_{ij}^{(k-1)}$

    \item \textbf{Il vertice $\mathbf{k}$ è un vertice intermedio del cammino}\\
    Il cammino esiste se e solo se esistono entrambi i sottocammini da $i$ a $k$ e da $k$ a $j$ con vertici intermedi in $1,\dots,k-1$, espresso dal termine $t_{ik}^{(k-1)} \land t_{kj}^{(k-1)}$
\end{itemize}

\section{Soluzione del problema principale}
Quando $k=n$, l'insieme dei possibili vertici intermedi coincide con l'intero insieme $V$.

La matrice finale della chiusura transitiva corrisponde a $T^{(n)} = (t_{ij}^{(n)})$

\section{Algoritmo Bottom-up}
\subsection{Algoritmo bottom-up}

\begin{algorithm}
    \begin{algorithmic}
        \Procedure{CHIUSURA-TRANSITIVA}{$M,n$}
            T = M \Comment{Caso base $k=0$: inzializzazione con la matrice di adiacenza}

            \For{$k=1$ to $n$}
                \For{$j=1$ to $n$}
                    \State $t_{i,j}=t_{i,j} \lor (t_{i,k} \land t_{k,j})$
                \EndFor
            \EndFor

            \State \Return $T$
        \EndProcedure
    \end{algorithmic}
\end{algorithm}

\subsection{Complessità}
\begin{itemize}
    \item \textbf{Complessità temporale:} $\Theta(n^3)$, dovuta ai tre cicli for annidati su $k$, $i$ e $j$ che eseguono operazioni logiche elementari a costo costante $\Theta(1)$.
    \item \textbf{Complessità spaziale:} $\Theta(n^2)$
\end{itemize}
